The three-epoch matrix was followed by a separate extension of the
observed-first conditions in both lines. PU-GCN, the strict point budgets
and the observation-preservation rule were retained. PointRCNN completed
18 RPN epochs in each line, followed by 18 RCNN epochs for Line~A and a
24-epoch RCNN schedule for Line~B. CenterPoint completed 12 epochs per
line. Each epoch checkpoint was evaluated on the same 3,769 validation
frames. The PointRCNN loader retained 3,265 effective training frames
after object and range filtering; 3,712 denotes the source training
split, not the number of frames remaining after that filtering.

Validation progress was assessed using Car Moderate 3D AP$_{R40}$.
Starting with the first epoch's score, a running reference was updated
only when a subsequent score exceeded it by more than 0.2 AP points.
Such an update reset the counter of consecutive epochs without a
threshold-exceeding improvement; otherwise the counter increased by
one. A completed schedule met the plateau criterion when this counter
was at least three at its end. All epochs were inspected, including any
later improvement after an intermediate plateau. The reported checkpoint
was the numerical AP maximum over the schedule, which need not coincide
with the last update of the threshold reference. The 0.2 threshold is
an operational training rule, not a statistical significance test.

Changes in the total number of epochs were treated as new schedules.
In particular, the 24-epoch Line~B RCNN experiment restarted from the
same RCNN initialisation with the selected RPN epoch~14 fixed. Its
one-cycle learning-rate trajectory therefore differed from the earlier
18-epoch run; the two trajectories were not concatenated. Results from
the earlier 12-epoch RPN and 18-epoch Line~B RCNN schedules remain
historical comparisons rather than additional segments of the final
curves.

The extension did not retrain every baseline for the same duration.
Line~A retained the original-$N$ input with official detector weights
as its reference. Line~B retained the sparse-$M$ baseline after the
earlier three-epoch adaptation, with the frozen sparse baseline also
reported. The original-$N$ reference is additionally shown to make the
remaining deficit relative to the complete observation explicit.
Checkpoint selection and final reporting used the same validation set.
Consequently, these are validation-selected system comparisons with
unequal training budgets, not independent test results or an isolated
causal effect of point generation. The plateau criterion applies to the
selected Car metric; it does not establish convergence of every class
or of the optimisation parameters.
